% TikZ Figure Definitions for FedRouteAI Manuscript with Auto-Scaling Bounds
\definecolor{ieeeblue}{RGB}{0,102,204}
\definecolor{lightbg}{RGB}{245,247,250}
\definecolor{bordergray}{RGB}{120,120,120}

\tikzset{
    module/.style={draw=black!80, line width=0.5pt, fill=white, rounded corners=3pt, align=center, font=\sffamily\fontsize{8}{9.5}\selectfont, inner sep=4pt},
    accentmodule/.style={draw=ieeeblue, line width=0.8pt, fill=ieeeblue!5, rounded corners=3pt, align=center, font=\sffamily\fontsize{8}{9.5}\selectfont, inner sep=4pt},
    subsystem/.style={draw=bordergray, line width=0.5pt, dashed, fill=lightbg!50, rounded corners=4pt, inner sep=6pt},
    arrow/.style={-{-Latex[length=4pt, width=3pt]}, line width=0.6pt, draw=black!80},
    accentarrow/.style={-{-Latex[length=4pt, width=3pt]}, line width=0.8pt, draw=ieeeblue},
    lbl/.style={font=\sffamily\fontsize{7}{8}\selectfont, text=black!70}
}

% Figure 1: Overall FedRouteAI System Architecture
\newcommand{\drawFigSystemArchitecture}{
\begin{figure*}[t]
\centering
\resizebox{0.95\textwidth}{!}{%
\begin{tikzpicture}[node distance=1.2cm and 1.5cm]
    % Edge Layer
    \node[module] (client1) {Edge Node 1\\Telemetry};
    \node[module, right=0.6cm of client1] (client2) {Edge Node 2\\Telemetry};
    \node[module, right=0.6cm of client2] (clientN) {Edge Node K\\Telemetry};
    \node[subsystem, fit={(client1) (client2) (clientN)}] (edge_sub) {};
    \node[above=0.1cm of edge_sub.north, font=\sffamily\bfseries\fontsize{8}{9.5}\selectfont] {Edge Client Layer (Decentralized Nodes)};

    % Memory Layer
    \node[accentmodule, below=1.2cm of edge_sub] (memory) {Historical Memory Module\\(PostgreSQL Database -- Zone Statistics $v_{avg}, d_{avg}$)};
    \node[module, right=1.2cm of memory] (feat) {Feature Engineering\\Normalized Ratio Vectors ($v/v_{hist}, d/d_{hist}$)};

    % Inference Engine
    \node[accentmodule, right=1.2cm of feat] (mlp) {EventClassifierMLP\\3,240 Params (Edge Inference)};

    % Server & Aggregator
    \node[module, above=1.2cm of mlp] (server) {Central Backend Server\\(FastAPI Router \& Flower FL Server)};
    \node[accentmodule, right=1.2cm of server] (agg) {AdaptiveFedAvg Aggregator\\Weight Scaling ($w_k = n_k q_k \rho_k g_k$)};

    % Dashboard
    \node[module, below=1.2cm of agg] (dash) {React Dashboard\\Real-Time Anomaly Alerting};

    % Flow Arrows
    \draw[arrow] (client1) -- (memory);
    \draw[arrow] (client2) -- (memory);
    \draw[arrow] (clientN) -- (memory);
    \draw[accentarrow] (memory) -- node[lbl, above] {Historical Baselines} (feat);
    \draw[accentarrow] (feat) -- node[lbl, above] {Relative Features} (mlp);
    \draw[arrow] (mlp) -- node[lbl, left] {12.96 KB Weights} (server);
    \draw[accentarrow] (server) -- node[lbl, above] {Client Updates} (agg);
    \draw[arrow] (agg) |- node[lbl, right] {Updated Global Model} (mlp);
    \draw[arrow] (mlp) -- node[lbl, above] {Predictions} (dash);
\end{tikzpicture}%
}
\caption{Overall FedRouteAI system architecture illustrating the decoupled interaction between decentralized edge clients, the PostgreSQL Historical Memory Module, the lightweight EventClassifierMLP inference engine, and the server-side AdaptiveFedAvg aggregator.}
\label{fig:architecture}
\end{figure*}
}

% Figure 2: Historical Memory Module Dataflow
\newcommand{\drawFigHistoricalMemory}{
\begin{figure}[htbp]
\centering
\resizebox{\linewidth}{!}{%
\begin{tikzpicture}[node distance=1.0cm and 1.2cm]
    \node[module] (telemetry) {Incoming Telemetry\\(GPS Speed $v$, Density $d$, Zone ID)};
    \node[accentmodule, below=0.8cm of telemetry] (db) {PostgreSQL Historical Store\\(Spatial-Temporal Zone Averages)};
    \node[module, below=0.8cm of db] (calc) {Relative Ratio Engine\\$v_{ratio} = v / v_{hist}$, $d_{ratio} = d / d_{hist}$};
    \node[accentmodule, below=0.8cm of calc] (tensor) {Normalized Feature Tensor\\(10-Dimensional Vector)};

    \draw[arrow] (telemetry) -- node[lbl, right] {Raw Metrics} (db);
    \draw[accentarrow] (db) -- node[lbl, right] {Zone Baselines} (calc);
    \draw[accentarrow] (calc) -- node[lbl, right] {Standardized Inputs} (tensor);
\end{tikzpicture}%
}
\caption{Historical Memory module dataflow depicting real-time queries against spatial-temporal baseline statistics ($v_{avg}, d_{avg}$) to compute normalized relative telemetry ratios.}
\label{fig:historical_memory}
\end{figure}
}

% Figure 3: Feature Engineering Pipeline
\newcommand{\drawFigFeatureEngineering}{
\begin{figure}[htbp]
\centering
\resizebox{\linewidth}{!}{%
\begin{tikzpicture}[node distance=0.8cm and 1.0cm]
    \node[module] (raw) {Raw Vehicle Telemetry\\$[v, d, lat, lon, t]$};
    \node[module, right=0.8cm of raw] (norm) {Zone Relative Scaling\\$v / v_{hist}, d / d_{hist}$};
    \node[accentmodule, right=0.8cm of norm] (vec) {10D Feature Vector\\$\mathbf{x} \in \mathbb{R}^{10}$};

    \draw[arrow] (raw) -- (norm);
    \draw[accentarrow] (norm) -- (vec);
\end{tikzpicture}%
}
\caption{Feature engineering pipeline transforming raw absolute GPS speed and traffic density telemetry into standardized 10-dimensional relative ratio vectors.}
\label{fig:feature_engineering}
\end{figure}
}

% Figure 4: Deployment-Aware Inference Pipeline
\newcommand{\drawFigInferencePipeline}{
\begin{figure}[htbp]
\centering
\resizebox{\linewidth}{!}{%
\begin{tikzpicture}[node distance=0.8cm and 1.0cm]
    \node[module] (input) {Raw Telemetry Vector};
    \node[module, below=0.7cm of input] (impute) {Deterministic Zero-Imputation\\(Missing Attribute Fallback)};
    \node[accentmodule, below=0.7cm of impute] (ood) {OOD Threshold Check\\($v_{ratio} > 3.0$ or $d_{ratio} > 7.0$)};
    \node[module, left=1.0cm of ood] (warn) {Generate Dynamic\\OOD Warning Flag};
    \node[accentmodule, below=0.7cm of ood] (mlp) {EventClassifierMLP\\Inference (0.2069 ms)};
    \node[module, below=0.7cm of mlp] (out) {Anomaly Prediction + Warning Header};

    \draw[arrow] (input) -- (impute);
    \draw[arrow] (impute) -- (ood);
    \draw[arrow] (ood) -- node[lbl, above] {Threshold Exceeded} (warn);
    \draw[accentarrow] (ood) -- node[lbl, right] {Valid / Warning Flagged} (mlp);
    \draw[accentarrow] (mlp) -- (out);
    \draw[arrow] (warn) |- (out);
\end{tikzpicture}%
}
\caption{Deployment-aware inference execution flow incorporating deterministic zero-imputation for missing attributes and dynamic Out-of-Distribution (OOD) threshold alert generation ($speed\_ratio > 3.0$).}
\label{fig:inference_pipeline}
\end{figure}
}

% Figure 5: AdaptiveFedAvg Aggregation Workflow
\newcommand{\drawFigAdaptiveFedAvg}{
\begin{figure}[htbp]
\centering
\resizebox{\linewidth}{!}{%
\begin{tikzpicture}[node distance=0.8cm and 1.0cm]
    \node[module] (updates) {Client Updates ($w_t^k, L_k, n_k$)};
    \node[accentmodule, below=0.7cm of updates] (quality) {Quality Factor: $q_k = e^{-L_k}$};
    \node[accentmodule, right=0.8cm of quality] (rarity) {Rarity Bonus: $\rho_k = 1/\sqrt{F_{normal}^k}$};
    \node[accentmodule, below=0.7cm of quality] (gate) {Reliability Gate: $g_k = 1 - \alpha \max(0, V_k - V_{thresh})$};
    \node[module, below=0.7cm of rarity] (weight) {Weighting: $\omega_k = n_k \cdot q_k \cdot \rho_k \cdot g_k$};
    \node[accentmodule, below=0.7cm of weight] (norm) {Normalize: $\hat{\omega}_k = \omega_k / \sum \omega_j$};
    \node[module, below=0.7cm of norm] (agg) {Global Model: $w_{t+1} = \sum \hat{\omega}_k w_t^k$};

    \draw[arrow] (updates) -- (quality);
    \draw[arrow] (updates) -| (rarity);
    \draw[arrow] (quality) -- (gate);
    \draw[arrow] (gate) |- (weight);
    \draw[arrow] (rarity) -- (weight);
    \draw[accentarrow] (weight) -- (norm);
    \draw[accentarrow] (norm) -- (agg);
\end{tikzpicture}%
}
\caption{Server-side AdaptiveFedAvg aggregation workflow showing the multi-factor dynamic weight calculation combining dataset volume ($n_k$), validation quality ($q_k$), rarity bonus ($\rho_k$), and temporal reliability gating ($g_k$).}
\label{fig:adaptivefedavg_workflow}
\end{figure}
}

% Figure 6: Federated Learning Communication Workflow
\newcommand{\drawFigFLCommunication}{
\begin{figure*}[t]
\centering
\resizebox{0.95\textwidth}{!}{%
\begin{tikzpicture}[node distance=1.0cm and 2.0cm]
    \node[accentmodule] (server) {Flower FL Server\\Global Model $w_t$};
    \node[module, below left=1.2cm and 2.5cm of server] (c1) {Client 1\\Local Train (5 Epochs)};
    \node[module, below left=1.2cm and 0.5cm of server] (c2) {Client 2\\Local Train (5 Epochs)};
    \node[module, below right=1.2cm and 0.5cm of server] (c3) {Client 3\\Local Train (5 Epochs)};
    \node[module, below right=1.2cm and 2.5cm of server] (cK) {Client 8\\Local Train (5 Epochs)};

    \draw[accentarrow] (server) -- node[lbl, left] {1. Broadcast $w_t$ (12.96 KB)} (c1);
    \draw[accentarrow] (server) -- (c2);
    \draw[accentarrow] (server) -- (c3);
    \draw[accentarrow] (server) -- node[lbl, right] {1. Broadcast $w_t$} (cK);

    \draw[arrow] (c1) -- node[lbl, right] {2. Upload $w_{t+1}^1, L_1$} (server);
    \draw[arrow] (c2) -- (server);
    \draw[arrow] (c3) -- (server);
    \draw[arrow] (cK) -- node[lbl, left] {2. Upload $w_{t+1}^8, L_8$} (server);
\end{tikzpicture}%
}
\caption{Federated learning communication protocol detailing parameter exchanges (12.96 KB per client) across 8 simulated edge nodes.}
\label{fig:fl_communication}
\end{figure*}
}

% Figure 10: Complete End-to-End Experimental Pipeline
\newcommand{\drawFigExperimentalPipeline}{
\begin{figure*}[t]
\centering
\resizebox{0.95\textwidth}{!}{%
\begin{tikzpicture}[node distance=1.0cm and 1.2cm]
    \node[module] (gen) {Synthetic Telemetry\\Dataset (30,195 Samples)};
    \node[module, right=0.8cm of gen] (part) {Non-IID Data Partition\\(8 Edge Clients)};
    \node[accentmodule, right=0.8cm of part] (fl) {Flower FL Simulation\\(30 Communication Rounds)};
    \node[accentmodule, right=0.8cm of fl] (eval) {Held-Out Evaluation\\(Centralized vs. AdaptiveFedAvg)};
    \node[module, right=0.8cm of eval] (bench) {Latency \& Communication\\Benchmarking (0.2069 ms)};

    \draw[arrow] (gen) -- (part);
    \draw[arrow] (part) -- (fl);
    \draw[accentarrow] (fl) -- (eval);
    \draw[accentarrow] (eval) -- (bench);
\end{tikzpicture}%
}
\caption{End-to-end experimental evaluation pipeline from non-IID data partitioning across 8 clients through federated training to systems-level latency benchmarking.}
\label{fig:experimental_pipeline}
\end{figure*}
}
